\section{Invasion of the Body Snatchers}

\begin{quotex}
The reason that people don't take the red pill is not necessarily because they are afraid to face reality, but rather because, if you do, it becomes really annoying to hang around blue pill people. \flright{\textsc{moi}}

To begin with unlimited liberty is to end with unlimited despotism. \flright{\textsc{Fyodor Dostoevsky}}

The wise person seeks what he lacks, and he considers not how much he knows, but of how much he is ignorant. \flright{\textsc{Hugh of St. Victor}}

\end{quotex}
There have been several versions of the film, \emph{Invasion of the Body Snatchers}, most recently starring Nicole Kidman. The general idea is that aliens have infected the minds of human beings. They then become robotic, thinking and acting alike. A few who manage to avoid infection find themselves bewildered and confused. They are considered to be enemies of the state by those how are infected. They are pressured to conform to the others, or they at least have to pretend to conform.

If you accept the idea of the Revelation of the Method, then the film is not depicting some possible future, but is actually describing the present state of affairs. Some unknown ``they", in reality guided by demons, are able to alter attitudes and values of millions of people. Those unable or unwilling to conform then face intense pressure to do so. Of course, ``they" reveal to us what they are doing by means of the mass media. Not that it makes any difference, since few people have a Real I, and must necessarily follow the latest fads.

\paragraph{The Idea of the Good}
Plato noted that we are guided by two principles: our inborn desire for pleasure and the desire for what is objectively best. The first desire is based on serotonin: if it feels good, it must be good. Of course, following the crowd, i.e., being on the ``right side of history" might also feel good. The ``pleasure principle" is the state of fallen man at birth, so it does not need to be learned. This is the unregenerate level of fallen human beings. As such, it is contrary to nature, although most consider it to be natural.

That is not the case for the preference of the good over the pleasurable. That must be learned and requires the rational intellect to determine and to will the good. \textbf{David Hume}, in the \emph{Treatise of Human Nature} totally distorts this distinction when he asserts:

\begin{quotex}
Reason is and ought to be the slave of the passions.

\end{quotex}
That is the complete opposite of Tradition which claims that the rational faculty ought to rule over the passions. Only then can a human being exist at a higher level, which is his true natural state. Such a person has a Real I.

\paragraph{The Classes of Men}
There are three karmic classes of men:

\begin{itemize}
\item \textbf{Third class}. These are the unregenerate people. They are blindly seeking the meaning of life, without a real inner direction. They think they have an I who speaks and acts, but that I is not permanent; it changes from day to day, or even from minute to minute. They are ``tossed to and fro and carried about with every wind of doctrine by the sleight of men and their cunning and craftiness."They falsely believe they have a real I because all their subpersonalities respond to the same name and occupy the same body. Unfortunately, since their ideas do not come from their rational soul, they fall prey to the body snatchers. In extreme cases, as in multiple personality disorders, the subpersonalities have different names and often do not even know of the others. Those of the third class are not really different, although they may function better in the world. 
\item \textbf{Second class}. People in this class have become dissatisfied with the world of those of the third class. They then seek the meaning of life in science, philosophy, religion, etc., in short, in systems outside themselves. They may be quite intelligent and erudite, having accumulated a large amount of book knowledge. They may even know a lot of what is written on esoteric and Hermetic topics. However, they don't bother with the exercises so they never advance.They console themselves with the claim that the teachings are for a ``future generation". Hence, their lives, in practice, are not much different from those of the third class. 
\item \textbf{First class}. This class is composed of the few who have developed a sense of the real I. They understand that spiritual growth requires self-discipline and intentional suffering. It is essential to oppose one's faults, negative emotions, idle fantasies, etc., as we have often noted. Just reading about such practices is not at all effective. Such men will try to help others on the path, even if few respond. 
\end{itemize}
\paragraph{The Three Forces or Gunas}
The traditional teaching is that there are three forces behind events in the world. In Hindu metaphysics, these are the three gunas:

\begin{itemize}
\item \textbf{Rajas}. The force to change, i.e., the positive force. 
\item \textbf{Tamas}. The negative force that retards. 
\item \textbf{Sattva}. Harmonises the positive and negative forces. 
\end{itemize}
\textbf{Empedocles} came close to this understanding with his notions of Love and Strife, with Love as the attractive, or positive, force and Strife as the repulsive, or negative, force. Unfortunately, it seems that he was unable to discern the neutralizing or equilibrating force. It is always difficult to see the neutralizing force in the world. That is because the lower intellect is bipolar or dual, so conventional thinking finds it difficult to reconcile those opposites.

\textbf{Boris Mouravieff} called these forces the static, dynamic, and neutralizing principles. He describes them this way:

\begin{quotex}
Where the third or neutralizing force is concerned, it often escapes our observation—either due to the bipolar character of our minds, or because of its own nature—which can in many cases leave it unobserved. It also sometimes plays the role of catalyst —far less conspicuous than the binding tie which is the most fundamental role of the neutralizing force.

The passive force contains all the possibilities for creating the phenomenon, while the active force intervenes as the realizer, and the neutralizing force as the regulator, of the relations between the two other forces, determining the dosage for both in an optimal way.

\end{quotex}
Mouravieff gives several examples of how these forces operate in different circumstances. The ultimate example concerns the relationship of God with the World. God is the universal I and the world is the universal Thou. Love, then, is the neutralizing force between them, as described in John's Gospel. That is the solution to the problem of the One and the Many.

\paragraph{The Lotus Flowers}
Valentin Tomberg takes a different approach in his discussion of the chakras or lotus flowers. The Heart chakra, which is the midpoint, is that of Love. The three chakras above the Heart relate to the Spirit and the three lower chakras to the body, or world. The Heart chakra, as Love, then balances the two forces. The Heart chakra lives according to solar law; it then harmonizes the unruly lower chakras to also follow the solar law. It will also bring warmth to the abstract intellectuality of the higher chakras.

\paragraph{Love, Fear, and Hubris}
In the 1980s, three interesting texts appeared: \emph{Gnosis} by Mouravieff, \emph{Meditations on the Tarot} by Tomberg, and the allegedly channelled text, \emph{A Course in Miracles} [ACIM]. I immersed myself in all three. The text of ACIM was accompanied by 365 daily exercises. I followed them for one year; I had a box of 365 cards and carried the appropriate one in order to remember to do the exercise. Obviously, doing such exercises has had a lasting effect on me. Primarily, the idea of doing a task at specified times during the day has stuck with me. It forces one to pause and question his habitual way of understanding the world. The text did not interest me as much, as it has some problematic elements; hence, I no longer have any interest in it. At the time, I knew several people involved with the course. I suspect, however, that few of them actually did the exercises.

The point I'm getting to is that I have some credentials to comment on Marianne Williamson, who, as a presidential candidate, has intrigued and confounded many. It seems to me to be difficult to reconcile ACIM with partisan politics. Specifically, the Course teaches Love which should obliterate dualistic thinking. Yet, like Empedocles, she divides the world into Fear and Love. Obviously, she claims to be motivated by Love, while Trump feeds on fomenting Fear. This is what the Course says:

\begin{quotex}
Only love is strong because it is undivided. The strong do not attack because they see no need to do so. Before the idea of attack can enter your mind, you must have perceived yourself as weak. … The opposite of seeing through the body's eyes is the vision of Christ, which reflects strength rather than weakness, unity rather than separation, and love rather than fear.

\end{quotex}
Gnosis is based on spiritual vision, not on our sensory perception of the World:

\begin{quotex}
Misperceptions produce fear and true perceptions foster love, but neither brings certainty because all perception varies. That is why it is not knowledge. True perception is the basis for knowledge, but knowing is the affirmation of truth and beyond all perceptions.

\end{quotex}
Gnosis begins with self-observation and self-knowledge. Such knowledge is always certain and beyond doubt:

\begin{quotex}
The Bible tells you to know yourself, or to be certain. Certainty is always of God. When you love someone, you have perceived him as he is, and this makes it possible for you to know him.

\end{quotex}
This applies, a fortiori, to our love for God. We only love what we know. When we love God, we love his commandments, which are simply descriptions of the Solar Law. To follow the Solar Law, then, is to serve God.

Ultimately, Miss Williamson has made the Course irrelevant. On the stage, she supported the exact same policies as all the others. This worldview is independent of any personal beliefs. We saw that Jews, atheists, liberal Christians, secularists, and so on, all came to the same conclusion. The impulse for this must be arising from a much different, and unnoticed, force. It's as though an alien force has snatched their bodies;

Since they have no explanation for their political viewpoint, they resort to confabulations, after the fact, to explain it. These confabulations lead to hubris, since they see themselves as spiritually enlightened, saviours of humanity, superior to their opponents, etc.

\paragraph{Unconditional Love}
Many women have told me that their dogs have taught them unconditional love. That is because dogs are poor judges of character: Hitler's dog loved Hitler, Stalin's dog loved Stalin, and Pol Pot's dog loved Pol Pot. A dog will love whoever feeds it and gives it attention. The dog's attitude is quite natural.

As an experiment in unconditional love, go to your neighbour's house and demand a free meal. Before you leave, drop a deuce on her carpet. If she is still showing you unconditional love, ask her to rub your belly before you go.

The Christian teaching is that we are commanded to love, so a dog is incapable of love. On the face of it, that sounds quite odd, since we consider love to be passive not active; e.g., one ``falls in love". Everyone assumes he knows what love is, but it is actually not so easy to understand. Even when understood, it is difficult to love, and the spiritual authorities have not been all that helpful.

\paragraph{God's Kingdom}
Although it is true that God is Love, there are many misunderstandings about what that means in practice. But before there can be Love, there must be fear; fear of the Lord is the beginning Wisdom.

God rules, not just by love, but by these three factors.

\begin{itemize}
\item Fear of Hell 
\item Promise of Heaven 
\item God rules by love: persuasion and creativity 
\end{itemize}
For most people, at least most religious people, the fear of Hell and the promise of Heaven are the prime motivating factors. To love God is to follow the commandments for their own sake, and not because of Hell or Heaven.

That is the path to self-mastery, which expels the body snatchers.



\flrightit{Posted on 2019-07-18 by Cologero }
